% Generated draft: requires mathematical and source review.
\documentclass[11pt]{article}
\usepackage[margin=25mm]{geometry}
\usepackage{fontspec}
\setmainfont{Latin Modern Roman}
\usepackage{amsmath,amssymb,mathtools,unicode-math}
\setmathfont{Latin Modern Math}
\usepackage{enumitem,needspace,xurl,hyperref}
\setlength{\parindent}{0pt}
\setlength{\parskip}{5pt}
\setlength{\emergencystretch}{3em}
\newcommand{\N}{\mathbb{N}}
\newcommand{\Z}{\mathbb{Z}}
\newcommand{\Q}{\mathbb{Q}}
\newcommand{\R}{\mathbb{R}}
\newcommand{\C}{\mathbb{C}}
\newcommand{\F}{\mathbb{F}}
\newcommand{\A}{\mathbb{A}}
\newcommand{\PP}{\mathbb{P}}
\newcommand{\E}{\mathbb{E}}
\newcommand{\eps}{\varepsilon}
\newcommand{\dd}{\,\mathrm{d}}
\begin{document}
\section*{UnsolvedMath ID: 700700}
\textbf{Source problem code:} catalogue entry
\subsection{Definitions and mathematical statement}
Let \(G\) be a group. Define its lower central series by
\[
\gamma_1(G)=G,\qquad \gamma_{n+1}(G)=[\gamma_n(G),G]\quad(n\geq 1),
\]
where \([H,G]\) is the subgroup generated by all commutators \(h^{-1}g^{-1}hg\) with \(h\in H\) and \(g\in G\). Write
\[
(*)\qquad \exists n\geq 1\;\;\gamma_n(G)\text{ is a free group}.
\]
Here a free group means a group having a free basis, with no restriction on the cardinality of that basis. Under the usual convention, the trivial group is free on the empty basis. The source does not specify a different convention or exclude \(n=1\).

A group \(G\) is residually finite if, for every \(g\in G\setminus\{1\}\), there exist a finite group \(F\) and a homomorphism \(\varphi:G\to F\) such that \(\varphi(g)\neq 1\).

\begin{enumerate}
\item[(a)] For a group given by a presentation \(G=\langle X\mid R\rangle\), find conditions on that presentation under which
\[
\bigl(\exists n\geq 1:\gamma_n(G)\text{ is free}\bigr)
\Longrightarrow G\text{ is residually finite}.
\]
No finiteness assumption on \(X\), \(R\), or \(G\) is supplied.
\item[(b)] Find conditions on the structure of a residually finite group \(G\) that ensure
\[
\exists n\geq 1:\gamma_n(G)\text{ is free}.
\]
\end{enumerate}
These are requests to identify hypotheses ensuring the indicated implications. The source does not specify a class of permitted hypotheses or require necessary and sufficient conditions.
\subsection{Short English statement}
Find presentation conditions ensuring that a group with a free term in its lower central series is residually finite. Conversely, find structural conditions ensuring that a residually finite group has such a free term. The scope of the requested conditions is unspecified.
\subsection{Sources}
\begin{itemize}
\item Local UnsolvedMath registry, supplied complete record: permanent ID \(700700\), source problem code \texttt{catalogue entry}, title \emph{Kourovka Notebook Problem 11.7}; upstream revision \texttt{d371ee917ac9647fbb06cd6fc529deec6f3dda82}.
\item Supplied attribution: K. Bencs\'ath, proposed in 1990; \emph{The Kourovka Notebook (Unsolved Problems in Group Theory)}, original problem 11.7, 11th issue (1990). The supplied background also identifies the 21st issue edition (2026), \texttt{arXiv:1401.0300v48}. The collection description names E. I. Khukhro and V. D. Mazurov as editors.
\item The local record labels this problem \emph{open} and records an update date of 2026-10-16. This is supplied registry information, not an independently verified assessment of mathematical or literature status.
\end{itemize}
\end{document}
